% Level 3 answer key, Week 12 -- Showcase and Graduation
\answersweek{12}{Showcase and Graduation}{What good work looks like for workbook chapter 12 and homework 12. Exit assessment answers are in the assessments section.}

\section{Exercise 12.1: Showcase talk card}
\begin{answer}[Good work looks like]
One real decision, not a feature list. The \emph{Why} gives a reason connected to a concept: \enquote{inheritance so I did not repeat code}, \enquote{a dictionary so I can look up scores by nickname}, \enquote{BFS because it finds the shortest path}, \enquote{a test caught a bug}. \emph{Next time} is honest and specific. The talk fits in 2 minutes.
\end{answer}

\section{Exercise 12.2: Feedback cards}
\begin{answer}[Good work looks like]
Stars and wishes are specific and kind: \enquote{The ghost going around walls was clever}, \enquote{What if the score was bigger?}. Not: \enquote{cool}, \enquote{it is bad}.
\end{answer}

\section{Exercise 12.3: My Level 3 journey}
\begin{answer}[Good work looks like]
The commit count matches \texttt{git log --oneline} (target 8+). Proudest skill and \enquote{hard then, easier now} are personal -- any honest answer is right. Use them in your end-of-level note to parents.
\end{answer}

\section{Homework 12}
\begin{answer}[Core: letter to future-you]
Look for the four parts: what they built, a bug they fixed, advice for hard moments, a goal for next year. Do not score spelling; this is personal. Suggest sealing it in an envelope with a date.
\end{answer}
\begin{answer}[Challenge and Creative]
Challenge: one menu item with a concrete, small first step (\enquote{add a second ghost type and commit}), planned with a parent. Creative: badge shows a nickname and three real skills.
\end{answer}

\begin{teachernote}[Exit assessment]
Answers and the rubric for the Level 3 exit assessment are in \enquote{Answers: assessments} at the end of this answer key.
\end{teachernote}
